\section{Distribución y streaming}

Una vez producida la señal de programa, esta debe distribuirse hasta la audiencia.
En producción en directo sobre IP, la distribución se realiza mediante protocolos de
streaming que equilibran de distinta forma tres factores en tensión: latencia,
fiabilidad y escalabilidad.

\subsection{Protocolos de streaming}

\begin{itemize}
  \item \textbf{RTMP} (\textit{Real-Time Messaging Protocol}): protocolo desarrollado
  por Adobe sobre TCP, con latencia típica de 2--5~s. Durante más de una década fue
  el estándar \textit{de facto} para ingestar vídeo en plataformas como YouTube Live
  o Twitch. Aunque obsoleto en el lado del reproductor (Flash), sigue siendo el
  protocolo de ingesta más soportado por encoders y mezcladores de software, incluido
  Voctomix mediante FFmpeg.
  \item \textbf{HLS} (\textit{HTTP Live Streaming}): protocolo de Apple basado en la
  división del flujo en segmentos de fichero servidos por HTTP. Su latencia habitual
  es de 10--30~s (modo estándar) o 2--5~s (Low-Latency HLS). Es el protocolo de
  distribución más compatible con reproductores y CDN actuales.
  \item \textbf{SRT} (\textit{Secure Reliable Transport}): protocolo de baja latencia
  (< 1~s) con corrección de errores ARQ, diseñado para contribución y distribución
  sobre redes no confiables. Está ganando adopción rápida como sustituto de RTMP en
  flujos de ingesta.
  \item \textbf{WebRTC} (\textit{Web Real-Time Communication}): estándar del W3C para
  comunicación multimedia en tiempo real directamente desde el navegador, con latencias
  de 100--500~ms. Su uso en distribución broadcast es limitado por la complejidad de
  escalar a grandes audiencias, pero es ideal para monitorización remota y retorno
  de comunicaciones.
\end{itemize}

\subsection{Plataformas y CDN}

Las plataformas de streaming (YouTube Live, Twitch, Facebook Live) actúan como
intermediarios que reciben la señal del encoder mediante RTMP o SRT, la recodifican
en múltiples calidades y la distribuyen a la audiencia a través de sus propias redes
de distribución de contenidos (CDN). Esto descarga al sistema de producción de la
responsabilidad de escalar la distribución, que puede involucrar decenas de miles de
espectadores simultáneos.

En instalaciones propias ---como la configuración de tres escenarios implementada en
este trabajo--- la distribución puede realizarse directamente sin plataforma
intermediaria, empleando un servidor HLS o SRT accesible en la red local o con
acceso controlado desde Internet.

\subsection{Latencia frente a calidad}

Existe un compromiso fundamental entre latencia y calidad en los sistemas de
streaming: los codecs de alta eficiencia (H.265, AV1) requieren mayor complejidad
de codificación y ventanas temporales más largas para explotar la correlación entre
fotogramas, lo que se traduce inevitablemente en mayor latencia. En producción en
directo, la elección del punto de operación en este compromiso depende del caso de uso:
una retransmisión deportiva interactiva exige baja latencia aunque ello suponga menor
calidad, mientras que una conferencia técnica puede asumir latencias de decenas de
segundos a cambio de una imagen de mayor resolución y nitidez.
